% ============================================================
\section{Conclusion}
\label{sec:conclusion}
% ============================================================

We investigated whether the quality of synthetic negatives affects
scale robustness when fine-tuning foundation models for IoT
anomaly detection.  Three key findings emerge.

First, \textbf{learned negatives maintain higher detection across
training scales}: D-DiffTS matches Gaussian noise at 200/5
(AUC$=$0.621 vs.\ 0.629), improves at 500/10 (0.680 vs.\ 0.592),
and stabilises at 2000/30 (0.578 vs.\ 0.506; full-range
degradation $-$6.9\% vs.\ $-$19.6\%).
With a frozen encoder, learned negatives achieve the best AUC in
the paper (0.728, $+$0.146 over frozen Gaussian).
A generator ablation isolates this advantage to the learned base
distribution rather than post-hoc guidance.
Second, \textbf{catastrophic forgetting} drives scale degradation
for all negative types: freezing the encoder recovers AUC
(frozen B$=$0.606, C$=$0.582, D$=$0.573),
and converging evidence from a learning-rate ablation and LoRA experiment
rules out hyperparameter misconfiguration.
Third, the NLL$+$SupCon training recipe transfers between the two
architectures tested: it improves detection for both Moirai and from-scratch CNN, and
generalises cross-domain to N-BaIoT with Gaussian negatives
(AUC$=$0.928$\pm$0.041); D-DiffTS on N-BaIoT with 310 training windows
achieves only AUC$=$0.660$\pm$0.427, but with 1{,}096 windows (device~4)
reaches AUC$=$0.939$\pm$0.014, confirming that learned negatives require
$\geq$500 samples and match Gaussian when adequate data is available.

The practical implication: when fine-tuning foundation models for
anomaly detection, invest in generative-model-quality training
negatives rather than increasing training scale with low-quality
negatives.  When generative training is not feasible, freeze the encoder;
use CNN at small scale; always recalibrate thresholds on
deployment-specific benign traffic.
We caution that no configuration achieves operationally useful F1 at
stealth-95 (best: 0.262); the system is a diagnostic framework for
understanding detection failure modes, not a production-ready IDS.
We release all code, datasets, and calibration scripts.
